% !TEX encoding = UTF-8 Unicode
% !TEX program = pdflatex



\documentclass[border=1cm, tikz]{standalone}

\def\lightx@{.3}
\def\lighty@{-.5}
\pgfdeclarefunctionalshading[sphere-color]{sphere}
	{\pgfpoint{-25bp}{-25bp}}{\pgfpoint{25bp}{25bp}}
	{\pgfshadecolortorgb{sphere-color}{\spherecolor}}{
	25 div exch 25 div exch 2 copy 
	dup mul exch dup mul add 1 sub
	\lightx@ dup mul \lighty@ dup mul add 1 sub
	mul abs sqrt
	exch \lightx@ mul add
	exch \lighty@ mul add
	.6 mul .3 add
	dup abs add 2 div
	\spherecolor
	3 index mul 3 1 roll
	3 index mul 3 1 roll
	3 index mul 3 1 roll
}% https://tex.stackexchange.com/questions/54193/how-to-draw-a-shaded-sphere/75994#75994

\begin{document}
\begin{tikzpicture}
	\begin{scope}% inside black
		\fill (0,0) circle [radius=4.9];
	\end{scope}
	\begin{scope}% upper red
		\clip (-5,0.3) -- (-5,5) -- (5,5) --(5,0.3)
			arc[start angle=0, end angle=-180, x radius=5cm, y radius=1.5cm];
		\colorlet{sphere-color}{red}
		\shade[shading=sphere] (0,0) circle[radius=5cm];
	\end{scope}
	\begin{scope}% lower white
		\clip (-5,-0.3) -- (-5,-5) -- (5,-5) --(5,-0.3)
			arc[start angle=0, end angle=-180, x radius=5cm, y radius=1.5cm];
		\colorlet{sphere-color}{white}
		\shade[shading=sphere] (0,0) circle[radius=5cm];
	\end{scope}
	\begin{scope}% button
		\fill[black] (-2,-1.4) circle [x radius=1.5cm, y radius=1.2cm, rotate=-45];
		\fill[white] (-2,-1.4) circle [x radius=1cm, y radius=0.8cm, rotate=-45];
	\end{scope}
\end{tikzpicture}
\end{document}